\section{题04：安全路径（P16786）}

\subsection{问题建模：子树奇偶性与有序点对计数}

题目给出一棵以 $1$ 为根的 $n$ 个节点的无向树。
记 $\mathit{sz}[x]$ 为以 $x$ 为根的子树节点总数（含 $x$ 自身）。
若 $\mathit{sz}[x]$ 为偶数，称 $x$ 为平稳节点。
对于不同节点 $x,y$，若树中唯一简单路径上每个节点都平稳，
则有序对 $(x,y)$ 是一条安全路径，$(x,y)$ 与 $(y,x)$ 视为两条不同路径。
输入 $n$ 和 $n-1$ 条边，输出安全路径总数，$n \le 500000$，
要求线性或近线性时间算法。
本质是在树上按局部性质筛选节点子集，再统计子集内满足全平稳条件的点对数量。

\subsection{暴力枚举的局限：$\mathrm{O}(n^2)$ 在 $n=5\times10^5$ 下的困境}

最朴素思路：$\mathrm{O}(n)$ 算出 $\mathit{sz}$ 和平稳性，
再枚举 $\mathrm{O}(n^2)$ 个有序点对，逐对验证路径节点平稳性。
即便用 LCA 优化路径提取，单次验证最坏扫 $\mathrm{O}(n)$ 个节点，
总复杂度 $\mathrm{O}(n^3)$。即使能做到 $\mathrm{O}(1)$ 回答任意点对是否安全，
$n^2 \approx 2.5\times10^{11}$ 次查询本身已严重超时。
$n=5\times10^5$ 下，必须放弃显式枚举，寻找结构性计数方法。

\subsection{核心洞察：平稳节点诱导子图是森林}

重新审视条件：若 $x$ 到 $y$ 路径全平稳，则路径上相邻节点对也必平稳。
整条路径完全位于平稳节点诱导子图中——仅保留平稳节点，
且只保留两端均平稳的边。树连通无环，诱导子图必为森林。

在此森林中，$x$ 与 $y$ 之间有安全路径当且仅当同属一个连通分量：
原树两点简单路径唯一，全平稳则留在诱导子图中，两点连通；
连通路径也必全平稳。因此，大小为 $s$ 的连通分量内，
不同两点均构成两条有序安全路径，贡献 $s\cdot(s-1)$。
答案是所有分量贡献之和。

以示例验证：$n=6$，边 $(1,2),(1,3),(1,6),(2,4),(3,5)$，根 $1$。
$\mathit{sz}[4]=\mathit{sz}[5]=\mathit{sz}[6]=1$（奇，非平稳），
$\mathit{sz}[2]=\mathit{sz}[4]+1=2$（偶，平稳），
$\mathit{sz}[3]=\mathit{sz}[5]+1=2$（偶，平稳），
$\mathit{sz}[1]=\mathit{sz}[2]+\mathit{sz}[3]+\mathit{sz}[6]+1=6$（偶，平稳）。
平稳集 $\{1,2,3\}$。两端均平稳的边 $(1,2),(1,3)$ 形成大小 $3$ 的分量，
贡献 $3\times2=6$，与示例吻合。

\subsection{算法推导：自底向上DP中的分量合并与屏障结算}

可在诱导子图上 DFS 得分量大小，但显式构造需额外筛边。
更优雅的做法是原树上自底向上 DP，边算分量边累加答案，
全程仅对树两次遍历。

第一步：BFS 得层级与子树大小。从根 BFS，得 $\mathit{parent}[u]$ 和层次序
$\mathit{order}$（根到叶）。逆序处理：将 $\mathit{sz}[u]$ 累至 $\mathit{sz}[\mathit{parent}[u]]$，
完成后 $\mathit{stable}[u]=(\mathit{sz}[u]\bmod2=0)$。

第二步：自底向上 DP 算分量。再次逆序处理节点 $u$，
令 $\mathit{comp}[u]$ 为 $u$ 子树内以 $u$ 为顶的平稳连通分量大小。
若 $u$ 非平稳，$u$ 成屏障——向祖先的路径被 $u$ 中断。
遍历子 $v$：若 $v$ 平稳，分量无法再上延，结算
$\mathit{ans}\gets\mathit{ans}+\mathit{comp}[v]\cdot(\mathit{comp}[v]-1)$。
若 $u$ 平稳，$\mathit{cur}=1$，遍历子 $v$：
$v$ 平稳则边 $(u,v)$ 两端均平稳，合并 $\mathit{cur}\gets\mathit{cur}+\mathit{comp}[v]$；
$v$ 非平稳则其分量已在 $v$ 处结算，跳过。$\mathit{comp}[u]=\mathit{cur}$。

全节点处理完后，若根 $1$ 平稳，其分量未被子孙屏障结算，
手动追加 $\mathit{ans}\gets\mathit{ans}+\mathit{comp}[1]\cdot(\mathit{comp}[1]-1)$。

以链状树演示：$n=4$，边 $(1,2),(2,3),(3,4)$，根 $1$。
$\mathit{sz}[4]=1$（奇），$\mathit{sz}[3]=2$（偶），$\mathit{sz}[2]=3$（奇），$\mathit{sz}[1]=4$（偶）。
平稳 $\{1,3\}$，被非平稳 $2$ 隔开。
BFS 序 $[1,2,3,4]$，逆序：
$4$ 叶且非平稳，无操作；
$3$ 平稳，$\mathit{cur}=1$，子 $4$ 非平稳，$\mathit{comp}[3]=1$；
$2$ 屏障，子 $3$ 平稳，$\mathit{ans}+=1\times0=0$；
$1$ 平稳，$\mathit{cur}=1$，子 $2$ 非平稳，$\mathit{comp}[1]=1$；
结算根 $\mathit{ans}+=1\times0=0$。答案为 $0$，符合预期。

\subsection{复杂度与实现：线性扫描、long long 与迭代遍历}

两次遍历：BFS 建序 $\mathrm{O}(n)$，逆序累加 $\mathit{sz}$ 和 DP 各 $\mathrm{O}(n)$，
每条边各访问一次，总时间 $\mathrm{O}(n)$。空间 $\mathrm{O}(n)$（邻接表加五个线性数组），
$n=500000$、256 MB 下绰绰有余。

实现要点。答案须用 64 位有符号整数（C++ \mil{long long}）：
最坏全树平稳且连通，$n(n-1)\approx2.5\times10^{11}$ 远超 32 位 \mil{int} 范围。
树深可达 $n$（退化链），递归 DFS 会栈溢出，必用迭代 BFS。
$\mathit{order}$ 逆序即自底向上序。邻接表 \mil{vector<vector<int>>}，
读边双向建边，BFS 用 $\mathit{parent}$ 防回访。
结算公式 $s(s-1)$ 以 64 位运算：\mil{1LL * comp[v] * (comp[v] - 1)}。

\subsection{代码实现}
\inputminted{c++}{../src/p04_safe_paths.cpp}
